% ==============================================================================
% CAPÍTULO 2: IMPLEMENTAÇÃO PRÁTICA DOS CLIENTES IOT E DO BACKEND DE IRRIGAÇÃO
% ==============================================================================

\begin{center}
\setstretch{1.0}
\textbf{\fontsize{14}{16}\selectfont CAPÍTULO 2}

\vspace{0.5em}

\textbf{\fontsize{14}{16}\selectfont IMPLEMENTAÇÃO PRÁTICA DOS CLIENTES IOT E DO BACKEND DE IRRIGAÇÃO}
\end{center}

\vspace{1.5em}

\section*{1. INTRODUÇÃO À IMPLEMENTAÇÃO DE CLIENTES E BACKEND EM IRRIGAÇÃO IOT}

A consolidação de um sistema automatizado de irrigação depende da tradução prática das especificações teóricas do protocolo HTTP em código-fonte executável e funcional. Enquanto o Capítulo 1 estabeleceu a semântica das requisições, respostas, verbos e códigos de status HTTP, este capítulo dedica-se exclusivamente à construção da camada de implementação dos múltiplos clientes e do servidor central backend.

No ecossistema de automação agrícola proposto, a heterogeneidade de dispositivos exige abordagens de programação distintas para cada perfil de cliente. Os nós embarcados de campo operam sob restrições de hardware e memória RAM, demandando código em linguagem C++ otimizado para os microcontroladores ESP32 e Arduino. No ambiente do escritório ou central de monitoramento, o painel de controle exige reatividade gráfica orientada a eventos em TypeScript com o framework Angular. Por fim, o servidor backend necessita de uma estrutura robusta em Python Flask integrada ao banco de dados relacional MySQL para gerenciar a persistência e a tomada de decisão em tempo real.

Nas seções a seguir, são apresentadas e explicadas as implementações completas dos clientes de hardware embarcado, da aplicação web frontend e da API do backend de irrigação (\texttt{backend-irriga}).

\newpage

\section*{2. IMPLEMENTAÇÃO DOS CLIENTES EMBARCADOS (EDGE HARDWARE)}

Os dispositivos de campo são responsáveis por realizar a interface direta com o mundo físico, acionando relés elétricos que operam a bomba de água principal e as válvulas solenoides dos setores de cultivo.

\subsection*{2.1. Código C++ para ESP32: Sincronização via HTTPClient}

O microcontrolador ESP32 executa um firmware em C++ que estabelece a conexão Wi-Fi com a rede local da fazenda e efetua requisições periódicas do tipo \texttt{GET /estado-geral} para o servidor Flask. O código analisa a resposta JSON recebida, extrai o estado booleano do atributo \texttt{bomba} e altera o nível lógico do pino digital conectado ao módulo relé.

\begin{lstlisting}[language=C++]
#include <WiFi.h>
#include <HTTPClient.h>
#include <ArduinoJson.h>

const char* ssid = "Rede_Fazenda_IoT";
const char* password = "Senha@Irrigacao2026";
const char* serverUrl = "http://192.168.1.100:5000/estado-geral";

const int PIN_RELE_BOMBA = 26;

void setup() {
  Serial.begin(115200);
  pinMode(PIN_RELE_BOMBA, OUTPUT);
  digitalWrite(PIN_RELE_BOMBA, LOW);

  WiFi.begin(ssid, password);
  while (WiFi.status() != WL_CONNECTED) {
    delay(500);
    Serial.print(".");
  }
  Serial.println("\nWi-Fi Conectado!");
}

void loop() {
  if (WiFi.status() == WL_CONNECTED) {
    HTTPClient http;
    http.begin(serverUrl);
    http.addHeader("Accept", "application/json");

    int httpCode = http.GET();

    if (httpCode == HTTP_CODE_OK) {
      String payload = http.getString();
      StaticJsonDocument<512> doc;
      DeserializationError error = deserializeJson(doc, payload);

      if (!error) {
        bool bombaAtiva = doc["bomba"];
        digitalWrite(PIN_RELE_BOMBA, bombaAtiva ? HIGH : LOW);
        Serial.printf("Estado da Bomba Atualizado: %s\n", bombaAtiva ? "LIGADA" : "DESLIGADA");
      }
    } else {
      Serial.printf("Erro na Requisicao HTTP: %d\n", httpCode);
    }
    http.end();
  }
  delay(5000);
}
\end{lstlisting}

A análise técnica do código demonstra a eficiência da biblioteca \texttt{HTTPClient}, que abstrai a abertura e o fechamento do socket TCP, permitindo que a decodificação da estrutura JSON ocorra via memória estática alocada pela biblioteca \texttt{ArduinoJson}.

\newpage

\subsection*{2.2. Código C++ para Arduino UNO com Ethernet Shield W5100}

Em cenários rurais onde a infraestrutura sem fio sofre com interferências eletromagnéticas de motores elétricos, utiliza-se a plataforma Arduino UNO acoplada a um módulo Ethernet cabeado W5100. Devido às limitações extremas de memória RAM do microcontrolador ATmega328P (2 KB de RAM), a requisição é transmitida via socket TCP bruto através do objeto \texttt{EthernetClient}.

\begin{lstlisting}[language=C++]
#include <SPI.h>
#include <Ethernet.h>

byte mac[] = { 0xDE, 0xAD, 0xBE, 0xEF, 0xFE, 0xED };
IPAddress server(192, 168, 1, 100);
int port = 5000;

EthernetClient client;

void setup() {
  Serial.begin(9600);
  if (Ethernet.begin(mac) == 0) {
    Serial.println("Falha ao configurar Ethernet via DHCP");
    for (;;);
  }
  delay(1000);
}

void loop() {
  Serial.println("Conectando ao Servidor de Irrigacao...");

  if (client.connect(server, port)) {
    // Construcao manual da requisicao HTTP GET RAW
    client.println("GET /bomba HTTP/1.1");
    client.println("Host: 192.168.1.100");
    client.println("User-Agent: ArduinoEthernet/1.0");
    client.println("Connection: close");
    client.println(); // Linha em branco obrigatoria

    // Aguarda e le a resposta do servidor
    while (client.connected() || client.available()) {
      if (client.available()) {
        char c = client.read();
        Serial.print(c);
      }
    }
    client.stop();
    Serial.println("\nConexao encerrada.");
  } else {
    Serial.println("Falha na conexao TCP com o servidor.");
  }
  delay(10000);
}
\end{lstlisting}

Nessa implementação de baixo nível, a transmissão manual das linhas de cabeçalho e da linha em branco de término (\texttt{CRLF}) evidencia o funcionamento bruto do protocolo HTTP/1.1 conforme especificado na RFC 9112 (FIELDING; NOTTINGHAM; RESCHKE, 2022b).

\newpage

\section*{3. IMPLEMENTAÇÃO DO CLIENTE WEB ANGULAR}

O painel de controle web permite ao agricultor visualizar graficamente o estado do sistema, cadastrar novas rotinas de irrigação e acionar manualmente válvulas solenoides. A aplicação é estruturada como uma Single Page Application (SPA) em Angular.

\subsection*{3.1. Serviço HTTP Reativo em TypeScript}

O serviço \texttt{IrrigacaoService} encapsula as chamadas à API REST do backend. Utilizando o módulo \texttt{HttpClient} do Angular, o serviço retorna objetos assíncronos do tipo \texttt{Observable} gerenciados pela biblioteca RxJS.

\begin{lstlisting}
import { Injectable } from '@angular/core';
import { HttpClient, HttpHeaders } from '@angular/common/http';
import { Observable } from 'rxjs';

export interface Agendamento {
  id?: number;
  data_agendamento: string;
  hora_inicio: string;
  hora_fim: string;
  duracao_minutos: number;
  solenoide: string;
  usuario: string;
  status: string;
}

@Injectable({
  providedIn: 'root'
})
export class IrrigacaoService {
  private apiUrl = 'http://192.168.1.100:5000';

  constructor(private http: HttpClient) {}

  // Obtem o estado unificado do sistema
  getEstadoGeral(): Observable<any> {
    return this.http.get<any>(`${this.apiUrl}/estado-geral`);
  }

  // Aciona bomba ou solenoide
  setSolenoide(id: number, ativo: boolean): Observable<any> {
    return this.http.post<any>(`${this.apiUrl}/solenoide?id=${id}&ativo=${ativo}`, {});
  }

  // Cadastra um novo agendamento de irrigacao
  criarAgendamento(agendamento: Agendamento): Observable<Agendamento> {
    const headers = new HttpHeaders({ 'Content-Type': 'application/json' });
    return this.http.post<Agendamento>(`${this.apiUrl}/agendamentos`, agendamento, { headers });
  }
}
\end{lstlisting}

A interface \texttt{Agendamento} define a tipagem estrita do modelo de dados, garantindo que requisições enviadas ao backend atendam aos requisitos de validação do contrato da API.

\newpage

\section*{4. IMPLEMENTAÇÃO DO BACKEND EM PYTHON FLASK (\texttt{backend-irriga})}

O servidor central backend é desenvolvido na linguagem Python com o micro-framework Flask e o ORM SQLAlchemy para integração com o banco de dados MySQL.

\subsection*{4.1. Lógica de Sincronização e Rotas REST API}

O módulo \texttt{bomba\_routes.py} implementa o roteamento das requisições HTTP e executa a função de sincronização automática entre a tabela de agendamentos e a tabela do atuador da bomba de água.

\begin{lstlisting}[language=Python]
from flask import Blueprint, jsonify, request
from sqlalchemy import text
from db import db

bomba_routes = Blueprint('bomba_routes', __name__)

def sincronizar_status():
    # 1. Verifica agendamentos ativos na data e horario atuais
    sql_check = """
    SELECT COUNT(*) FROM agendamentos 
    WHERE data_agendamento = CURDATE() 
      AND CURTIME() BETWEEN hora_inicio AND hora_fim
    """
    agendamentos_ativos = db.session.execute(text(sql_check)).scalar() or 0
    deve_ligar_bomba = agendamentos_ativos > 0

    # 2. Atualiza o status da bomba principal no banco de dados
    db.session.execute(
        text("UPDATE bomba SET ativo = :ativo WHERE id = 1"),
        {"ativo": 1 if deve_ligar_bomba else 0}
    )
    db.session.commit()

@bomba_routes.route('/estado-geral', methods=['GET'])
def estado_geral():
    try:
        sincronizar_status()
        solenoides = db.session.execute(text("SELECT id, ativo FROM solenoides")).fetchall()
        bomba = db.session.execute(text("SELECT ativo FROM bomba WHERE id = 1")).scalar()

        return jsonify({
            "bomba": bool(bomba),
            "solenoides": [{"id": s[0], "ativo": bool(s[1])} for s in solenoides]
        }), 200
    except Exception as e:
        return jsonify({"error": f"Erro interno: {str(e)}"}), 500
\end{lstlisting}

A função \texttt{sincronizar\_status()} avalia se a hora atual do servidor encontra-se compreendida dentro de algum intervalo de irrigação agendado no banco relacional. Caso existam registros ativos, o status da bomba é atualizado para 1 (\texttt{true}), assegurando que os nós embarcados recebam a instrução de acionamento durante a consulta periódica ao endpoint \texttt{/estado-geral}.

\vspace{1.5em}

\section*{REFERÊNCIAS}

\begin{singlespace}
\noindent FIELDING, Roy T.; NOTTINGHAM, Mark; RESCHKE, Julian. \textbf{HTTP Semantics}. RFC 9110, Internet Engineering Task Force (IETF), jun. 2022. Disponível em: <https://www.rfc-editor.org/rfc/rfc9110>. Acesso em: 28 set. 2026.

\vspace{0.8em}
\noindent FIELDING, Roy T.; NOTTINGHAM, Mark; RESCHKE, Julian. \textbf{HTTP/1.1}. RFC 9112, Internet Engineering Task Force (IETF), jun. 2022. Disponível em: <https://www.rfc-editor.org/rfc/rfc9112>. Acesso em: 28 set. 2026.

\vspace{0.8em}
\noindent KRUGER, C. P. \textbf{Automação Agrícola e Internet das Coisas}: Aplicação de Sensores e Microcontroladores no Campo. São Paulo: Editora Técnica, 2024.

\vspace{0.8em}
\noindent SANTOS, T. R.; SILVA, M. A. \textbf{Sistemas de Irrigação Inteligentes baseados em ESP32 e Protocolos Web}. Revista Brasileira de Engenharia Agrícola e Ambiental, v. 28, n. 4, p. 210-225, 2025.
\end{singlespace}
